Misurare, analizzare e classificare segnali fisiologici provenienti dal corpo umano rappresenta uno dei campi di ricerca più affascinanti e all'avanguardia del panorama scientifico. I progressi nella strumentazione, nelle tecnologie, nei paradigmi e negli algoritmi degli ultimi decenni hanno portato notevoli miglioramenti nell'assistenza all'uomo e alla sua cura. 

Nel campo della misurazione dell'attività cerebrale, ricerche sempre più avanzate forniscono descrizioni dettagliate sul funzionamento delle reti cerebrali, sulle relazioni tra funzioni e specifiche aree o determinati ritmi motori e sul comportamento dinamico durante i processi cognitivi. Da queste informazioni derivano gli strumenti fondamentali per la caratterizzazione di malattie e alterazioni neurologiche, quali epilessia, demenze, Parkinson, Alzheimer. Sulla base di queste informazioni è oggi inoltre possibile pianificare interventi medici al cervello \cite{zhang2024braincomputer}. 

In questo contesto è rilevante il campo della classificazione dell'intenzione cerebrale in corrispondenza di determinati ambienti, quali l'immaginazione motoria e la stimolazione esterna. Nata a partire dal 1971 \cite{Nirenberg1971-rp} e da allora in costante sviluppo, la classificazione dei segnali cerebrali ha permesso la nascita di tecnologie di nuova frontiera quali, prime fra tutte, le \emph{Brain-Computer Interface} (BCI). Questi dispositivi utilizzano le informazioni estratte dall'attività cerebrale per mettere in comunicazione il cervello di un soggetto con un dispositivo esterno in modo tale da controllarlo assecondando l'intenzione umana.

Si distinguono due direzioni progettuali che mirano ad obiettivi diversi: sistemi BCI a ciclo aperto oppure a ciclo chiuso. I sistemi a ciclo aperto o unidirezionali pongono come obiettivo la realizzazione di una strumentazione a sostegno del soggetto. I sistemi a ciclo chiuso prevedono invece un meccanismo di feedback essenziale per le operazioni di neuroriabilitazione. 

I sistemi a ciclo aperto si appoggiano su una serie altamente variegata di strumenti e algoritmi che ne rendono difficile la sintesi e che saranno difatti approfonditi nel prossimo capitolo. Seguendo la tendenza di maggior rilievo, il lavoro ambisce a contribuire nell'ambito dei sistemi BCI sostenuti da misurazioni elettroencefalografiche (EEG) di superficie. L'EEG è allo stato attuale il metodo più popolare: presenta costi relativamente ridotti, elevata portabilità e misura direttamente l'attività neuronale. 

Il lavoro ha portato allo sviluppo di due strumenti che sfruttassero tutti i dati forniti dal dataset di riferimento, dettagliatamente descritto nel capitolo 3. In particolare sono stati sviluppati:
\begin{itemize}
    \item \textbf{Un classificatore asincrono} basato su geometria riemanniana e valutato sulla segmentazione dei dati in finestre scorrevoli (\emph{sliding windows}). A sostegno del classificatore è stato implementato un algoritmo di selezione delle finestre atto ad aumentare l'accuratezza al costo di una minore percentuale di classificazioni accettate.
    \item \textbf{Un predittore di performance} che individua una correlazione tra la prominenza dei ritmi mu a riposo e l'accuratezza dello stesso soggetto.
\end{itemize}

Per quanto concerne il classificatore, le sue caratteristiche si discostano dagli approcci utilizzati da lavori simili: la maggioranza degli articoli che impiegano lo stesso dataset adottano il meccanismo di \emph{epoching} e restituiscono, dunque, stime generose. Il primo obiettivo è stato quello di valutare quanto un algoritmo di classificazione sia in grado di apprendere nelle condizioni meno favorevoli che un sistema BCI reale possa incontrare. In particolare, oltre alle limitazioni imposte dal dataset, sono state aggiunte alcune scelte di valutazione deliberatamente severe:
\begin{itemize}
    \item \textbf{una quantità ridotta di dati}: nelle classificazioni binarie l'addestramento è effettuato su due sole run, ossia circa quattro minuti di registrazione;
    \item \textbf{una validazione Leave-One-Run-Out (LORO)}: il modello viene sempre valutato su un'intera run mai vista in fase di addestramento, così che finestre sovrapposte appartenenti allo stesso trial non possano comparire contemporaneamente nell'insieme di training e in quello di test;
    \item \textbf{una segmentazione asincrona in finestre scorrevoli}: il classificatore decide su finestre di due secondi, con passo di mezzo secondo, lungo l'intera registrazione, come avverrebbe in un utilizzo reale del dispositivo.
\end{itemize}

Le accuratezze riportate in questa tesi non sono pertanto direttamente confrontabili con quelle della letteratura e vanno lette come una \textbf{stima conservativa} di quanto sia possibile ottenere. Se l'algoritmo riesce ad estrarre informazione utile in queste condizioni, è ragionevole attendersi che la prestazione possa soltanto migliorare con soggetti addestrati, sessioni multiple e un meccanismo di feedback.

I risultati ottenuti riflettono direttamente la capacità di ogni soggetto di emettere informazione utile ai fini della classificazione. Le accuratezze misurate esprimono quante finestre, in percentuale rispetto al totale, sono state classificate correttamente nei vari compiti di immaginazione motoria: pugno sinistro contro pugno destro, piedi contro pugni, riposo contro attivazione. Tale definizione è approssimativa e sarà spiegata dettagliatamente nei successivi paragrafi. Infine, essendo il sistema rivolto ad applicazioni real-time, le singole classificazioni sono state concatenate in un sistema a cascata che esegue la classificazione tra tutti i possibili stati.

Il secondo obiettivo, strettamente legato al primo, è comprendere quanto della variabilità osservata dipenda dall'algoritmo e quanto dal soggetto. A questo scopo le prestazioni sono state analizzate soggetto per soggetto, sia confrontando compiti di natura diversa sia isolando i soggetti migliori, e sono state messe in relazione con il predittore di performance, che a partire da un solo minuto di registrazione a riposo stima in anticipo la capacità del soggetto di controllare il sistema.

La tesi è organizzata come segue. Il capitolo 2 presenta le principali tecnologie per la misurazione dell'attività cerebrale, i paradigmi di acquisizione dei dati e gli algoritmi di estrazione delle feature impiegati nei sistemi BCI. Il capitolo 3 descrive il dataset utilizzato, il protocollo sperimentale con cui è stato raccolto, l'organizzazione delle registrazioni in run, trial e finestre e i limiti di cui tenere conto nell'interpretazione dei risultati. Il capitolo 4 illustra l'architettura sviluppata: il pre-processing, la segmentazione in finestre scorrevoli, l'estrazione delle feature riemanniane, la classificazione con la relativa validazione e il predittore di illiteracy. Il capitolo 5 riporta i risultati sperimentali delle classificazioni binarie, del sistema completo a cascata e del predittore, con particolare attenzione alla variabilità fra soggetti. Il capitolo 6, infine, discute i risultati ottenuti alla luce degli obiettivi posti, ne trae le conclusioni e individua i limiti del lavoro e i possibili sviluppi futuri.
